% !TeX root = booklet-03.tex
\chapter*{مرور تمرین‌ها و منابع مشخصه انتقال \lr{DC}}
\addcontentsline{toc}{chapter}{مرور تمرین‌ها و منابع مشخصه انتقال \lr{DC}}
\section*{تمرین \secnum{۲.۱۶} و محل استخراج پاسخ}
\label{sec:booklet-three-ex16}
صورت تمرین \secnum{۲.۱۶}، استخراج تحلیلی خروجی برحسب ورودی را در
دو ناحیه \lr{B} و \lr{D} با
\(\beta_n=\beta_p\) و \(V_{tn}=|V_{tp}|=V_t\) می‌خواهد
\cite[صفحه چاپی ۹۷]{weste}.
در بخش \ref{sec:inverter-five-regions}، مدار و بایاس‌ها،
فرض بار، برابری جریان‌ها، حل درجه دوم و انتخاب ریشه فیزیکی
به‌طور کامل استخراج شده‌اند. روابط
\eqref{eq:inverter-b-quadratic} تا
\eqref{eq:inverter-b-transfer} پاسخ \lr{B} و روابط
\eqref{eq:inverter-d-quadratic} تا
\eqref{eq:inverter-d-transfer} پاسخ \lr{D} را می‌دهند.
در این استخراج، فرض \(V_{DD}>2V_t\) برقرار است.
در \lr{B}، \lr{nMOS} در اشباع و \lr{pMOS} در
\lr{Linear} است و در \lr{D}، دو نقش عوض می‌شوند.
شرط‌های انتخاب ریشه و رفتار مرزی در همان بخش کنترل شده‌اند.
حل تمرین \secnum{۲.۱۷} در بخش \ref{sec:noise-ex17} با تکیه بر همین
روابط انجام می‌شود. شکل و پاسخ تشریحی تمرین \secnum{۲.۲۰} نیز در
بخش \ref{sec:pass-ex20} آمده است.

\section*{پوشش کتاب، اسلایدها و تمرین‌های دوره}
جدول \ref{tab:booklet-three-coverage} مسیر مطالعه هر موضوع
را نشان می‌دهد. کل بخش \secnum{۲.۵} کتاب، از صفحه چاپی
۸۷ تا آغاز ۹۳، در این جزوه پوشش داده شده است.
\begin{table}[!htbp]
\centering
\small
\caption{نقشه پوشش بخش \secnum{۲.۵} کتاب و اسلایدهای \lr{CMOS Inverter DC Characteristics}.}
\label{tab:booklet-three-coverage}
\begin{tabularx}{\linewidth}{c>{\raggedleft\arraybackslash}Xc}
\toprule
بخش کتاب & موضوع و محل در جزوه & اسلاید\\
\midrule
\secnum{۲.۵.۱} & بخش \ref{sec:inverter-dc}، مشخصه و ناحیه‌های \lr{Inverter} & ۳ تا ۷\\
\secnum{۲.۵.۲} & بخش \ref{sec:beta-ratio}، نسبت قدرت و \lr{Equivalent Inverter} & ۸\\
\secnum{۲.۵.۳} & بخش \ref{sec:noise-margin}، سطوح منطقی و نویز & ۹ و ۱۰\\
\secnum{۲.۵.۴} & بخش \ref{sec:pass-transistor}، افت آستانه و \lr{TG} & ۱۲ و ۱۳\\
مکمل زمانی & بخش \ref{sec:pass-dc-time}، تعریف ورودی \lr{Step} و \lr{Ramp} & ۱۱\\
تمرین \secnum{۲.۱۶} & بخش \ref{sec:inverter-five-regions} و مرور حاضر & ۱۴\\
تمرین \secnum{۲.۱۷} & بخش \ref{sec:noise-ex17}، استخراج و کنترل عددی & ۱۴\\
تمرین \secnum{۲.۲۰} & بخش \ref{sec:pass-ex20}، بازترسیم و تحلیل گره‌ها & ۱۴\\
\bottomrule
\end{tabularx}
\end{table}
اسلایدهای ۱ و ۲، عنوان و فهرست موضوعات‌اند. موضوع‌های علمی
اسلایدهای ۳ تا ۱۳ در جدول پوشش دارند. فهرست تمرین‌های
اسلاید ۱۴ به تمرین‌های \secnum{۲.۱۶}، \secnum{۲.۱۷} و \secnum{۲.۲۰} کتاب ارجاع می‌دهد.
شماره صفحه فایل کتاب برابر شماره چاپی به‌اضافه ۲۷ است.
شماره‌گذاری روابط و شکل‌های این جزوه مستقل از مرجع است.

